\begin{ExerciseBlock}
\ExerciseHeading{05}{Schwingung einer Scheibe}{free_oscillations_oscillating_disk}
Eine Schnur läuft schlupffrei über eine zylindrische Scheibe. Ihre Enden sind an zwei gleichen Federn mit unbekannter Federkonstante $D$ befestigt. Wird die Scheibe verdreht und losgelassen, schwingt sie. Ihre Masse ist $m=\qty{2}{\kilogram}$, ihr Trägheitsmoment $J=mR^2/2$.
\begin{center}
\CachedTikzFigure{content/resources/chapters/free_oscillations/tikz/disk_and_springs.tex}
\end{center}
\begin{enumerate}
\item Stellen Sie die ungedämpfte Bewegungsgleichung auf und geben Sie eine Formel für $T_0$ an.
\item Berechnen Sie $D$ für die gemessene Periodendauer $T_0=\qty{0.6}{\second}$.
\item Die Abklingkonstante beträgt $\delta=\qty{0.01}{\per\second}$. Wie groß ist die Amplitude nach 33 Schwingungen bei einer Anfangsamplitude von $\varphi_0=\ang{60}$? Verwenden Sie wegen der schwachen Dämpfung $T\approx T_0$.
\end{enumerate}
\end{ExerciseBlock}
\UseExerciseSolution{free_oscillations_oscillating_disk}
